% Рисунок с доски: вязкопластическая модель (предположительно Бингама —
% Шведова): пружина μ последовательно с блоком, где демпфер η параллелен
% элементу сухого трения.
\begin{figure}[htbp]
\centering
\begin{tikzpicture}[>={Latex[length=4pt]}, scale=0.9, transform shape]
  % схема
  \begin{scope}
    \draw[thick] (1.6,3.2) -- (1.6,3.0);
    \pic[rotate=90] at (1.6,2.1) {spring};
    \node[font=\footnotesize, anchor=west] at (2.05,2.1) {$\mu$};
    \draw[thick] (1.6,1.2) -- (1.6,1.0);

    \pic[rotate=90] at (0,0) {stvenant};
    \node[font=\footnotesize, anchor=east] at (-0.55,0) {$\tau_s$};
    \pic[rotate=90] at (3.2,0) {damper};
    \node[font=\footnotesize, anchor=west] at (3.45,0.5) {$\eta$};

    \draw[thick] (0,0.9) -- (0,1.0) -- (3.2,1.0) -- (3.2,0.9);
    \draw[thick] (0,-0.9) -- (0,-1.0) -- (3.2,-1.0) -- (3.2,-0.9);
    \draw[thick] (1.6,-1.0) -- (1.6,-1.6);
  \end{scope}

  % отклик при постоянной скорости деформации
  \begin{scope}[xshift=6.4cm, yshift=-1.6cm]
    \draw[->] (0,0) -- (5.4,0) node[anchor=north east] {$t$};
    \draw[->] (0,0) -- (0,4.2) node[anchor=south east] {$\tau$};
    \draw[dashed, gray] (0,2.2) -- (5.0,2.2);
    \node[anchor=east, font=\footnotesize] at (-0.05,2.2) {$\tau_s$};
    \draw[very thick, blue!60!black] (0,0) -- (1.3,2.2);
    \draw[very thick, blue!60!black, domain=1.3:5.0, samples=120, smooth]
      plot (\x, {2.2 + 1.3*(1-exp(-(\x-1.3)/0.9))});
    \draw[dashed, gray] (0,3.5) -- (5.0,3.5);
    \node[anchor=east, font=\footnotesize] at (-0.05,3.6)
      {$\tau_s + 2\eta\dt\gamma$};
    \node[font=\footnotesize, anchor=west] at (0.25,1.1) {$2\mu$};
  \end{scope}
\end{tikzpicture}
\caption{Вязкопластическая модель: пока $\abs{\tau^*} \le \tau_s$, течения нет и
работает пружина; за пределом скорость течения задаётся вязкостью,
$2\eta\,\dt\gamma_2 = \tau^* - \tau_s\sgn\tau^*$. Справа — отклик при постоянной
скорости деформации: напряжение выходит на уровень
$\tau_s + 2\eta\,\dt\gamma$.}
\label{fig:bingham}
\end{figure}
